\chapter*{ВВЕДЕНИЕ}
\addcontentsline{toc}{chapter}{ВВЕДЕНИЕ}

Современная электронная коммерция на маркетплейсах требует от продавцов высокоуровневой аналитической обработки больших объемов данных о продажах, остатках и финансовых показателях. Развитие технологий веб-разработки и методов анализа данных позволяет создавать специализированные программные комплексы, способные автоматически извлекать, структурировать и интерпретировать информацию из отчетов платформ, трансформируя её в стратегические бизнес – показатели~\cite{problemSellers}.

Актуальность данной выпускной квалификационной работы обусловлена критической зависимостью рентабельности бизнеса на маркетплейсах от оперативного управления товарными запасами, точного расчёта юнит-экономики с учётом специфических комиссий платформы и глубокого анализа динамики продаж, а также автоматизации бизнес-процессов. 

Дополнительная значимость работы заключается в её ориентации на практическое применение в качестве стартап-проекта. Разрабатываемая система направлена на решение актуальных задач продавцов маркетплейсов, связанных с анализом продаж, управлением товарными запасами и оценкой эффективности бизнеса. В отличие от учебного прототипа, проект изначально проектируется с учётом возможности дальнейшего внедрения, масштабирования и коммерциализации на рынке цифровых сервисов для электронной коммерции.

Цель работы – разработка системы и создание специализированного программного комплекса с использованием API методов для аналитики продаж и прогнозирования остатков товаров на маркетплейсе, обеспечивающего автоматизированный расчёт ключевых метрик платформы, а также реализацию механизмов аутентификации и авторизации пользователей для разграничения доступа к аналитическим данным и обеспечения безопасности работы с платформой.

Для выполнения поставленной цели нужно выполнить следующие задачи:

\begin{enumerate}
\item Изучение, создание и внедрение возможности регистрации и автоматизации пользователей.
\item Анализ специфики данных, бизнес-процессов и отчётности маркетплейса, получение данных API методом.
\item Проектирование и реализация модульной архитектуры веб-приложения для интерактивной аналитики данных.
\item Создание пользовательского интерфейса, обеспечивающего эффективную визуализацию сложных данных (синхронные таблицы, цветовое кодирование статусов ликвидности) и интуитивное взаимодействие с системой.
\end{enumerate}


